\documentclass[12pt]{article}
\usepackage[utf8]{inputenc}
\usepackage[T1]{fontenc}
\usepackage{hyperref}
\title{Wampur Language Resources for K-12 Curriculum Development}
\author{Compiled by Claude Sonnet 5 (claude-sonnet-5) for Andr\'as Kornai}
\date{}
\begin{document}
\maketitle
\section{Introduction}
Wampur (ISO 639-3 waz) is an Austronesian language of Papua New Guinea with about 530 speakers (2000 census). Sources disagree on location: Wikipedia places it in Madang Province, while Joshua Project places speakers on the Wanton River in Morobe Province. It should not be confused with Wampar. Speakers are reported to be trilingual in Wampur, Tok Pisin and Adzera. Little educational material was found.
\section{Primary Educational Resources}
No primers, textbooks or school materials in Wampur were found in the searches.
\section{Linguistic and Reference Materials}
\begin{itemize}
\item Wikipedia overview: \url{https://en.wikipedia.org/wiki/Wampur_language}
\item ANU Archives item with Wampur lexis (Ross and Tryon, 1984), a word list, and field notes and a grammar sketch (Pataki Schweizer, 1963): \url{https://archivescollection.anu.edu.au/index.php/wampur-language-papua-new-guinea-includes-wampur-lexis-m-ross-and-d-tryon-1984-wampur-word-list-1984-and-field-notes-from-k-pataki-schweizer-1963-wampur-grammar-sketch-wampur-grammar-outline-syntactic-check-list-wampur-gahuan-word-list}
\item Numeral data (MPI-EVA): \url{https://lingweb.eva.mpg.de/channumerals/Wampur.htm}
\end{itemize}
\section{Religious and Community Resources}
\begin{itemize}
\item Joshua Project language profile, which lists Bible translation as started: \url{https://joshuaproject.net/languages/waz}
\item Joshua Project people profile, which gives literacy of 43\% for men and 15\% for women: \url{https://joshuaproject.net/people_groups/15847}
\item Global Recordings Network: \url{https://globalrecordings.net/en/language/3150}
\end{itemize}
\section{Digital and Audio Resources}
\begin{itemize}
\item PARADISEC catalogue records: \url{https://catalog.paradisec.org.au/repository/SH2/ANUA520_3505} and \url{https://catalog.paradisec.org.au/repository/SH2/ANUA520_3508}
\end{itemize}
\section{Curriculum Development Context}
Wampur is used in village life across generations. Documentation is thin: a word list and a grammar sketch, with no pedagogical materials found. Any curriculum would need new work with the community, starting from the archived materials.
\end{document}